\subsection{Completing and correcting the polynomial}
\label{sec:symbol-correction}
\hypertarget{source-symbol-correction}{}

The lifted columns prescribe the initial part of a polynomial. Completing it at the shared endpoints introduces error terms whose orders must be estimated before the edge tests are combined.

Use \(x\) at an even start, \(v,y\) at an odd top, and let
\(H_i=p^{d_i}\), so \(H_iP_i=q_j\). The other factors are
formed from the integral polynomials
\[
 K_i=\sum_{a=0}^{k-1}x_a^{H_i}y_a,\qquad
 W_{ih}=v_hx_0^{H_i}-v_0x_h^{H_i}\quad(1\le h<k),\qquad
 D_i=\sum_a v_a y_a .
\]
Here \(K_i\) denotes a polynomial only in this paragraph and the
factor table below. The capacities elsewhere are those of
Section~\ref{sec:segment-window}.

\begin{lemma}[Initial symbols on an adapted line]
At the identity line \(w^{-1}S_1K_i\) and \(w^{-1}S_1W_{ih}\)
have leading terms \(x_0^{H_i}y_0\) and \(v_h x_0^{H_i}\),
respectively, with remainders divisible by \(w^{H_i-1}\).
Both unnormalized polynomials vanish modulo \(w\) on every
adapted line. For the same-position pairing one has
\[
 S_gD_i=wD_i.
\]

\end{lemma}
\begin{proof}
At the identity, ordinary vector and scalar covector coordinates
acquire a factor \(w\). Term by term,
\[
  S_1K_i\equiv w x_0^{H_i}y_0,\qquad
  S_1W_{ih}\equiv w v_hx_0^{H_i}\pmod {w^{H_i}} .
\]
At another adapted line the vector modulo \(w\) is proportional
to the Frobenius twist of the first inverse row. Raising the
start row to \(H_i\) gives the twist at the odd top, since the
residue frame has the cycle period. This kills the minor \(W_{ih}\).
The ordinary covector columns of the forward frame, with indices
other than zero, annihilate that row, killing \(K_i\). In the
same-position pairing the frame matrices cancel exactly, giving
the asserted formula for \(D_i\).
\end{proof}

At a tie of level \(j_i\) use \(D_i^{P_i}\) a total of \(j_i\)
times. Its normalized symbol is \(\sum_a(v_ay_a)^{P_i}\) by
Frobenius modulo \(w\). The last even pairing \(D_+^{j_+}\) has
exact normalized symbol \(D_+^{j_+}\).

\begin{lemma}[Raw polynomial tail at a distinguished line]
\label{lem:raw-tail-correction}
For \(1\le n\le c\) and a product generator \(b\) of
\(\mathcal S(n,o)\) of actual order \(o_b\le o\), there is a clean
polynomial \(F_0\) of that order for which
\[
 S_1 F_0-b\equiv\sum_{e=1}^{n-1}w^e b_e\pmod{w^nO[X]},
 \qquad b_e\in\mathcal S(n-e,o_b+e).
\]
The assertion uses the bounds and margins in
Section~\ref{sec:block-pairing}.
\end{lemma}

\begin{proof}
\proofstep{The missing exponents in the true weights.}
Writing \(\rho_{ih},\tau_{ih}\) for the vector and covector
digits, the remaining order-bearing factors are
\[
\begin{array}{c|c|c}
 \text{factor}&\text{multiplicity}&\text{order per factor}\\ \hline
 K_i^{P_i}&h_i\ (i<\sigma)&P_i\\
 W_{ih}&\rho_{ih}\ (h>0)&1\\
 K_i&\tau_{i0}&1 .
\end{array}
\]
Set \(h_i=0\) at ties when using this table.
The free digits \(v_0^{\rho_{i0}}\prod_{h>0}y_h^{\tau_{ih}}\)
and the endpoint covector monomial of degree \(b_+-j_+\) are
unchanged in the identity substitution and integral at all lines.

Choose one even start \(j_0\). Using the vector letters \(x_j\)
at all starts, put
\[
 C_h=\prod_j x_{j0}^{q_j-1}-\prod_j x_{jh}^{q_j-1},
 \qquad R=x_{j_0,0}\prod_{h=1}^{k-1} C_h ,
\]
where the subscripts denote the start and coordinate.
At the identity \(C_h\) has order zero and its symbol is the
first monomial. At another line write \((a_h)_h\) for the first
inverse row over the residue field and \(s_j\) for the start indices.
Each substituted \(x_{jh}\) reduces to \(x_{j0}a_h^{p^{s_j}}\).

For \(a_h\ne0\) we have
\[
 \prod_j(a_h^{p^{s_j}})^{q_j-1}=1,
\]
since the factors are
successive conjugate ratios around the cycle. If \(a_0=0\)
the initial factor in \(R\) vanishes. If \(a_0\ne0\) there is
some \(h>0\) with \(a_h\ne0\), since this is a different line.
Then \(C_h\) vanishes. Thus \(R^c\) supplies the extra \(c\)
orders at all competitors. In counting perturbations regard
each \(C_h\) as a separate factor repeated \(c\) times and
put the pure power at \(j_0\) with the exact monomial factors.

\paragraph{An endpoint shared by adjacent segments.}

An endpoint shared by two segments has only one degree to spend. The contributions from both adjacent paths must be included when that degree is restored.

The starting scalar degree spent at \(j\) is
\[
 U_j=q_j\Big(b_0+\sum_i h_i\Big)
 +\sum_i H_i\Big(\sum_{h>0}\rho_{ih}+\tau_{i0}\Big)
 + c\big((k-1)(q_j-1)+{\bf1}_{j=j_0}\big).
\]
Here \(b_0\) belongs to segment \(j\). The available degree
is \(a_{+,j-1}-o_{+,j-1}\), the decoration at the end of
the preceding segment, with cyclic indices. The first term is
at most \(a-E_0\) by the definition of \(b_0\). Each summand
in the second term is at most \(kq_j\), since each digit is
less than \(P_i\). The last term is at most \(ckQ\).

For the relaxed bounds of Section~\ref{sec:block-pairing} we have
\(o_{+,j-1}\le C_{\rm end}+c A_1\). The reserve
\[
 E_0=C_{E1}\left(a/p^2+b'/p+
     \mathcal L N_1a/(pk)+\mathcal L b'/k\right)
      + C_{E2} k^2 N_1 Q^2 c
\]
with both constants at least 17 covers
\(C_{\rm end}+cA_1+N_1 kQ+ckQ\).
Indeed the terms of \(C_{\rm end}\) have respective coefficients
at most \(5,3,4,2\) in the first parentheses for \(p\ge5\).
The other three terms satisfy
\[
 cA_1+N_1kQ+ckQ\le15kN_1Qc.
\]

Since \(a_{+,j-1}\ge a\), the unused scalar degree is nonnegative.
Fill it by the pure power of \(x_{j0}\).

\proofstep{The common degrees at the endpoint.}
The product \(F\) now has exactly the true multidegrees. This can
also be checked directly on its leading symbol. The columns give
the chain word with its radix substitutions and the odd powers
\(v_0^{P_i z_i}\). The factors in the table give the remaining
digits and the initial covector scalar powers. The tied and final
pairings give the powers of \(E_i,E_+\), and filling at all starts
gives exactly the final decorations. Every factor and its symbol
have identical letter degrees. The sum of the orders in the
columns counts the reduced vector degrees with their radix
multipliers. Restoring the levels and adding the normal digits and
the initial scalar covector powers gives
\[
 o_b=\sum_{\text{segments}}\left\{
 \sum_i\left[P_i u_i+\sum_{h>0}\rho_{ih}+\tau_{i0}\right]+o_+
 \right\},
\]
the order of the generator.

We have proved that \(F_0=w^{-o_b}F\)
is clean, has integral substitution at the chosen line, and that
its normalized symbol there is the generator modulo \(w\).
The extra \(c\) at another line is added to the full \(o_b\)
because the universal orders of all the factors are still available there.

\paragraph{The error of a product.}
Consider a variable factor of the construction, of order \(d\)
and leading symbol \(s\). Its normalized error is
\(w^{-d}S_1 f-s\). For a monomial of normal degree \(a\) its
coefficient is divisible by
\[
 w^{\max(1,a-d)}.
\]
Indeed
the symbol is homogeneous of normal degree \(d\), and the
diagonal substitution on the integral polynomial supplies
\(w^a\) on this monomial. The exact factors, including the final
pairing and its residual covector, have zero error. Every variable
factor has degree at most \(Q\) in each true letter, and degree
zero in the endpoint covector. The three degree bounds used here are
\[
 H_i P_i=q_j\le Q\quad(K_i^{P_i}),\qquad P_i\le Q\quad(D_i^{P_i}),\qquad q_j-1\le Q\quad(C_h).
\]


\proofstep{Terms indexed by the error factors.}
Expand the product of normalized factors, choosing errors in a
nonempty multiset of factors. Expand those errors in monomials.
A term modulo \(w^n\) may be written \(w^e\) times an integral
scalar times the product of the remaining symbols and one monomial,
where \(1\le e<n\). Take \(e\) to be the sum of the displayed
lower bounds on its individual coefficient orders, ignoring terms
with this sum at least \(n\). At most \(e\) factors have been
replaced. The monomial replacing them has their total
multidegrees and normal degree at most their total order plus
\(e\). The entire term consequently has normal degree at most
\(o_b+e\).

In each segment separate the remaining reduced column product
\(G'\), of degrees \(d_{vi},d_{yi},d_{v+}\), and the surviving
pairing powers \(E_i^{j'_i},E_+^{j_+}\). Here
\[
 0\le j'_i\le j_i
\]

at the ties, and both are zero in the initial list. In the
true letters all remaining pieces comprise a single monomial
\(X^A\). Before the replacements its exponents at an odd top were
\[
 (\rho_{i0}+P_i z_i,\ (\rho_{ih})_{h>0}),\qquad
 (\tau_{i0}+P_i h_i,\ (\tau_{ih})_{h>0}).
\]
At the endpoint these were the scalar vector decoration
\(a_+-o_+\) and the residual ordinary covector exponents
\(\gamma_h\). Write \(\Delta_h\) for the signed change in the exponent of coordinate \(h\) of a given true letter. The two directional bounds are
\[
 \sum_h(\Delta_h)_+\le Qe,\qquad \sum_h(-\Delta_h)_+\le Qe,\qquad q=Qe.
\]
For each replacement one removes its monomial
part, possibly only the pure scalar part of a column, and inserts
a monomial of degree at most \(Q\). If part of the factor was kept
separate in \(G'\) or the \(E_i\)'s, it is only the former loss
which becomes smaller. All endpoint covector exponents remain
exactly as before.

\proofstep{The additional order from deleted columns.}
There is a second effect of deleting columns. The largest
first-coordinate weight in a generic prefix frame is additive on
products of nonzero polynomials. Apply the generic change of
basis to the prefix and then its diagonal torus at the first
coordinate. Over the resulting coefficient domain, Laurent degree satisfies
\[
 \deg_T(FG)=\deg_T F+\deg_T G\qquad(FG\ne0).
\]
Indeed the product of the two leading coefficients is nonzero. Each deleted reduced column
has degree at most one in each of the at most \(N_1\) covectors.
Its largest weight is at least \(-N_1\), a lower bound even
for each of its monomials. Thus \(G'\) has every prefix cutoff
of the full word relaxed by at most \(N_1 e\).

\paragraph{Recovering a block from the error.}
Divide each odd-top exponent of \(A\) by \(P_i\), using its
remainders as the new digits. Write \(Z,Y\) for the quotients
in coordinate zero of the vector and covector, and
\(V_h,T_h\) for the ordinary quotients, at the top under
consideration. Write \(V=\sum_{h>0}V_h\) and \(V_+\)
for the total ordinary vector exponent at the endpoint of
radix one. The unchanged true degrees give new coarse totals
\[
 A'_i=j'_i+d_{vi}+Z+V,\qquad
 C'_i=j'_i+d_{yi}+Y+\sum_{h>0}T_h .
\]

In particular the new digit class is compatible and
\(A_i-A'_i\le k\), since two sums of \(k\) remainders
differ by less than \(kP_i\). We have \(V,V_+\le q\),
\(Y\le h_i+q\), and \(Z\ge z_i-q\). For example, each
positive unit of an ordinary vector quotient costs at least
one unit of positive variation from a digit less than \(P_i\).
The other assertions follow by the same division with remainder.

At an initial top set \(w'_i=d_{vi}+V\), \(h'_i=Y\).
Then \(w'_i\le w_i+q\), \(h'_i\le h_i+q\), and the two
new labels are bounded by \(A'_i,C'_i\), respectively. At a
tie set
\[
              o'_i=j'_i+d_{vi}+Y+V .
\]
Here \(j'_i+d_{vi}\le o_i\) and \(Y\le q\). In order to
have \(o'_i\le A'_i\) we need \(Y\le Z\). This uses the
excluded band quantitatively. For any compatible digits,
\(A_i\ge(11/12)b'/Q-k\), since \(b_i\ge(11/12)b'\).
Using \(b'\ge2k^3Q^4c\), \(\Delta\ge40/k\),
\(2N_1+3\le k\) and \(k\ge5\), we obtain
\[
 \Delta A_i/4\ge (55/3) k^2 Q^3 c-10
              \ge c(A_1+2Q+3k)+4 .
\]

\proofstep{The degree reserves for the reconstructed block.}
In the second inequality use \(A_1\le13k^2Q\),
\(Q,c\ge1\). The tied cutoff
\(o_i\le(1-\Delta/4)A_i+(c-n)A_1\) therefore gives
\(z_i=A_i-o_i\ge 1+2q\). Thus \[
 \begin{gathered}
 Y\le Z, \\ o'_i\le A'_i, \\ o'_i\le o_i+2q.
\end{gathered}
\]
Set the unused labels to zero. Finally put
\(o'_+=j_++d_{v+}+V_+\). It is bounded by \(a_+\)
by homogeneity and by \(o_++q\) by the reduced final degree.
The relaxed block inequalities now all hold at precision
\(n-e\). In particular the old tied upper bound loses at
most \(k+2q\) on changing \(A_i,o_i\), and the extra
allowance is \(eA_1\).

To realize the term in this block, multiply \(G'\) by all the
ordinary quotient monomials in the corresponding coarse letters,
including those at the endpoint. At an initial top the scalar
quotients already agree with the required decorations. At a
tie factor the scalar powers as
\[
 Z_i^{Z}S_i^{Y}
       = Z_i^{Z-Y}(E_i-v_i\cdot y_i)^Y ,
\]
now writing \(v_i,y_i\) for the ordinary coarse letters and
\(Z_i,S_i\) for their extra scalar coordinates.
Here
\[
 Z-Y=A'_i-o'_i.
\]

Binomial expansion gives new levels
\(j'_i+m_i\), \(0\le m_i\le Y\), and
\(Y-m_i\) new ordinary pairings. The identities for
\(A'_i,C'_i\) prove that both letter degrees at each such level
are exact and that its level is allowed. They also check the
degrees in the initial list. At the endpoint the scalar vector
power is exactly \(a_+-o'_+\), and the residual covector still
has degree \(b_+-j_+\).

\proofstep{The order of the reconstructed top term.}
An ordinary covector monomial cannot increase the top weight
of a prefix, and the inserted ordinary vector degrees cost at
most \(s q\) up to the odd tops. A new ordinary pairing costs
nothing unless the prefix stops between its two letters, where
it costs at most its power
\[
 Y-m_i\le q.
\]
Only one such pair
can straddle a prefix. If the endpoint vector is included, no
interior pair straddles the cut and the endpoint costs at most
\(q\). The total added prefix weight is therefore at most
\((s+1)q\), in addition to the deletion cost \(N_1e\).
Meanwhile each increment \(A_i-u_i\) in the capacity decreases
by at most \(k+2q\). The capacity can fall by at most
\(s(k+2q)\) at any prefix.

All three losses fit inside the
margin \(eA_1\), since
\[
    N_1e+(s+1)q+s(k+2q)\le13kN_1Qe=eA_1 .
\]
All costs here are integers, so taking floors after subtracting
the margin preserves the inequality. This proves membership in
the required integral chain lattice at the reduced precision.
The order of the new block equals the normal degree of the
term. Each occurrence of a coarse ordinary vector contributes its radix
and every pairing contributes that radix once. The digits supply
the remaining ordinary vector and scalar covector exponents.
Alternatively this follows by inserting the reconstructed labels
in the displayed formula for \(o_b\). The global monomial
decomposes by its odd tops and the endpoint of each segment,
each true letter being included exactly once. Thus an error term \(T_e\) of order \(e\) satisfies
\[
 w^{-e}T_e\in\mathcal S(n-e,o_b+e).
\]
This proves
\eqref{eq:raw-tail-required} and completes the simultaneous lift
by the induction in Section~\ref{sec:block-pairing}.

The constants in these inequalities are absolute under the stated local size bounds.
Neither the number of columns nor the number of competing
lines appears in the loss at one precision step.
\end{proof}
